% ============================================================
% SECCIÓN 0: Resumen Ejecutivo
% ============================================================
\chapter*{Resumen Ejecutivo}
\addcontentsline{toc}{chapter}{Resumen Ejecutivo}

\noindent
\textbf{Proyecto:} Demeter --- Plataforma IoT para Investigación Agrícola \\
\textbf{Autor:} Daniel Montero Fernández (\texttt{danmonfer@us.es}) \\
\textbf{Versión documentada:} 1.0 \\
\textbf{Estado:} Sistema operativo en banco; integración hidráulica y validación
de campo en curso \\
\textbf{Documento:} Documentación técnica completa de arquitectura, protocolo,
firmware, servicios y SDK

\vspace{0.6em}
\noindent\rule{\textwidth}{0.4pt}

% ─────────────────────────────────────────────────────────────
\section*{Propósito}

Demeter es una plataforma IoT diseñada para automatizar la recolección de variables
ambientales y el control de actuadores en entornos de investigación agrícola
(invernaderos, cámaras de cultivo). Resuelve el problema de la \textit{toma manual
de datos} y la \textit{intervención humana en el riego}, dos limitaciones que
condicionan tanto la calidad estadística de los experimentos como la disponibilidad
del personal investigador. El sistema combina tres elementos en un solo producto:
\textbf{red de sensores y actuadores} sobre microcontroladores ESP32, \textbf{servicio
de borde} ejecutado en una Raspberry Pi, \textbf{servidor central} con API REST y
WebSocket para visualización y control en tiempo real, y un \textbf{SDK científico de
Python} pensado para que investigadores puedan obtener informes de datos sin asistencia
del equipo de soporte.

% ─────────────────────────────────────────────────────────────
\section*{Arquitectura del Sistema}

El sistema se descompone en cuatro capas con responsabilidades estrictamente
separadas. La frontera entre capas se materializa mediante un \textbf{protocolo
binario propietario} (Demeter, descrito en el Capítulo~\ref{sec:protocolo}) entre
los nodos ESP32 y la Raspberry Pi, y mediante una \textbf{API REST + WebSocket}
entre la Raspberry Pi y el servidor central.

\begin{table}[H]
    \centering
    \caption{Capas del sistema Demeter y tecnologías asociadas}
    \begin{tabularx}{\textwidth}{lXl}
        \toprule
        \textbf{Capa}                       & \textbf{Función principal}                                                                       & \textbf{Tecnología} \\
        \midrule
        Campo (firmware)                    & Sensado, actuación y enrutamiento entre nodos.                                                     & ESP32 + ESP-NOW (C\texttt{++}) \\
        Borde (Raspberry Pi)                & Pasarela edge: UART$\leftrightarrow$WebSocket, despacho de comandos.                              & Python 3 + asyncio \\
        Servidor (backend)                  & API REST/WS, autenticación, persistencia, cola de cálculos pesados.                               & FastAPI + PostgreSQL + Redis \\
        Cliente (frontend / SDK)            & Dashboard interactivo y SDK de análisis científico para investigadores.                           & React (TS) + Python (\texttt{demeter\_sdk}) \\
        \bottomrule
    \end{tabularx}
    \label{tab:exec_arch}
\end{table}

\textbf{Decisiones arquitectónicas clave.} El proyecto adopta tres compromisos de
diseño que recorren toda la documentación: \textit{(i)} un protocolo binario
propietario en lugar de JSON sobre el enlace UART, motivado por el ancho de banda y la
latencia en redes ESP-NOW; \textit{(ii)} \textbf{Pydantic} como puente único entre la
representación binaria del protocolo y los modelos de aplicación, evitando duplicar
contratos; \textit{(iii)} \textbf{Docker dual} (despliegue en GCP y Raspberry Pi local)
para que el mismo \textit{stack} se ejecute en producción y en entorno doméstico sin
divergencia. Las decisiones se documentan en detalle en cada capítulo y se
formalizarán como ADRs antes del cierre v1.0.

% ─────────────────────────────────────────────────────────────
\section*{Alcance de la Versión 1.0}

El alcance funcional cubierto por esta versión es deliberadamente acotado para
garantizar un sistema cerrado y mantenible:

\begin{itemize}[leftmargin=*, label=--]
    \item Red de hasta 5 nodos ESP32 (sensores y actuadores) sobre ESP-NOW.
    \item Telemetría continua con periodicidad configurable y \textit{Deep Sleep} para nodos sensores.
    \item Control remoto de bombas y actuadores discretos a través del dashboard.
    \item Persistencia relacional (PostgreSQL) con LIMS básico: registro botánico, experimentos y trazabilidad.
    \item Análisis científico desde el SDK con caché Parquet local, indicadores agronómicos (VPD, GDD, ET$_0$) y exportación a Excel/CSV.
\end{itemize}

\textbf{Fuera de alcance v1.0} (registrado para evitar \textit{feature creep}): app
móvil, notificaciones push, capa de inferencia ML/IA, multi-tenancy, integraciones
con APIs meteorológicas y refactor estructural del código existente.

% ─────────────────────────────────────────────────────────────
\section*{Estado de Validación}

El proyecto sostiene un banco de \textbf{140 tests automatizados} distribuidos en
\textbf{18 archivos} sobre las cuatro capas (Capítulo~\ref{sec:testing}):

\begin{table}[H]
    \centering
    \caption{Validación automatizada por capa (verificada el 2026-05-07)}
    \begin{tabular}{lrrr}
        \toprule
        \textbf{Capa}      & \textbf{Archivos} & \textbf{Tests} & \textbf{Tipo predominante} \\
        \midrule
        Firmware (ESP32)   & 5                 & 32             & Unitario (Unity + mocks) \\
        Raspberry Pi       & 5                 & 37             & Mixto (asyncio mocks) \\
        Backend (FastAPI)  & 4                 & 16             & Integración (TestClient + SQLite) \\
        SDK (Python)       & 4                 & 52             & Unitario (funciones puras) \\
        \midrule
        Frontend           & 2                 & 3              & 3 unit. \\
        \midrule
        \textbf{Total}     & \textbf{20}       & \textbf{140}   & 117 unit. + 23 integ. \\
        \bottomrule
    \end{tabular}
    \label{tab:exec_tests}
\end{table}

La validación de hardware (caudal, consumo, ciclos térmicos de las bombas) estaba
prevista mediante \textbf{spikes} estructurados con hipótesis verificables, y el
criterio final de \texttt{DONE} para la versión 1.0 iba a ser \textbf{30 días de
operación continua sin intervención manual}. Conviene decirlo de entrada: \textbf{ni
el spike ni el periodo de operación llegaron a ejecutarse}. La versión 1.0 se entrega
verificada en software y validada en banco, pero \textbf{sin caracterizar en campo}
y sin haber actuado nunca sobre una bomba real. El Capítulo~\ref{sec:testing} detalla
qué se midió y qué no, y el Capítulo~\ref{sec:conclusiones} explica por qué.

% ─────────────────────────────────────────────────────────────
\section*{Cómo Leer Este Documento}

El documento está organizado en seis capítulos principales y tres anexos de diagramas:

\begin{itemize}[leftmargin=*, label=--]
    \item \textbf{Introducción} (visión, objetivos, alcance) y \textbf{Capítulo~\ref{sec:protocolo}}
          (protocolo binario): orientación general y contrato de comunicación.
    \item \textbf{Capítulos de arquitectura}: firmware ESP32, servicio Raspberry Pi
          (Capítulo~\ref{sec:raspberry}), backend (Capítulo~\ref{sec:backend}) y
          SDK (Capítulo~\ref{sec:sdk}).
    \item \textbf{Capítulo~\ref{sec:testing}}: estrategia de testing, inventario
          verificable de tests y métricas medidas.
    \item \textbf{Anexos}: diagramas de flujo del protocolo, arquitectura
          desplegada de la Raspberry Pi y arquitectura del backend.
\end{itemize}

\noindent\rule{\textwidth}{0.4pt}
\vspace{0.4em}

\noindent\textit{Este documento se mantiene sincronizado con el repositorio del
proyecto. La fecha de última verificación de los datos numéricos (tests, métricas,
inventario de archivos) se indica explícitamente en cada tabla.}
